
\documentclass{article}
\usepackage{colortbl}
\usepackage{makecell}
\usepackage{multirow}
\usepackage{supertabular}

\begin{document}

\newcounter{utterance}

\centering \large Interaction Transcript for game `bbh', experiment `salient\_translation\_error\_detection', episode 7 with reflexion{-}llama{-}3.1{-}8b{-}instruct.
\vspace{24pt}

{ \footnotesize  \setcounter{utterance}{1}
\setlength{\tabcolsep}{0pt}
\begin{supertabular}{c@{$\;$}|p{.15\linewidth}@{}p{.15\linewidth}p{.15\linewidth}p{.15\linewidth}p{.15\linewidth}p{.15\linewidth}}
    \# & \multicolumn{2}{c}{Player} && \multicolumn{2}{c}{Game Master} \\
    \hline

    \theutterance \stepcounter{utterance}  
    & & & \multicolumn{4}{p{0.6\linewidth}}{
        \cellcolor[rgb]{0.9,0.9,0.9}{
            \makecell[{{p{\linewidth}}}]{
                \texttt{\tiny{[P1$\langle$GM]}}
                \texttt{You have attempted this same task before and did not succeed. These are the plans you wrote after those attempts.} \\
\texttt{Trial \#1: New plan: I failed to identify the correct type of error in the translations. Upon reflection, I realize that I was too focused on the changes in the translations and didn't carefully read the options. In the future, I will carefully read the options and consider each type of error before making a decision. Specifically, I will make sure to pay attention to the changes in named entities, numerical values, modifiers or adjectives, negation or antonyms, facts, and dropped content. If I'm unsure, I will take a closer look at the source sentence and the translation to determine the correct type of error.} \\
\\ 
\\ 
\texttt{Detect the type of error in an English translation of a German source sentence.} \\
\\ 
\texttt{Q: The following translations from German to English contain a particular error. That error will be one of the following types: Named Entities: An entity (names, places, locations, etc.) is changed to a different entity. Numerical Values: Numerical values (ordinals or cardinals), dates, and/or units are changed. Modifiers or Adjectives: The modifiers and adjectives pertaining to a noun are changed. Negation or Antonyms: Introduce or remove a negation or change comparatives to their antonyms. Facts: Trivial factual errors not pertaining to the above classes are introduced in the translations. Dropped Content: A significant clause in the translation is removed. Please identify that error.  Source: In der Liste der Baudenkmale in Lenzen (Elbe) sind alle Baudenkmale der brandenburgischen Stadt Lenzen (Elbe) und ihrer Ortsteile aufgelistet.} \\
\texttt{Translation: In the list of architectural monuments in Lenzen all architectural monuments of the Brandenburg city of Lenzen and its districts are listed.} \\
\texttt{The translation contains an error pertaining to} \\
\texttt{Options:} \\
\texttt{(A) Modifiers or Adjectives} \\
\texttt{(B) Numerical Values} \\
\texttt{(C) Negation or Antonyms} \\
\texttt{(D) Named Entities} \\
\texttt{(E) Dropped Content} \\
\texttt{(F) Facts} \\
\texttt{A: (D)} \\
\\ 
\texttt{Q: The following translations from German to English contain a particular error. That error will be one of the following types: Named Entities: An entity (names, places, locations, etc.) is changed to a different entity. Numerical Values: Numerical values (ordinals or cardinals), dates, and/or units are changed. Modifiers or Adjectives: The modifiers and adjectives pertaining to a noun are changed. Negation or Antonyms: Introduce or remove a negation or change comparatives to their antonyms. Facts: Trivial factual errors not pertaining to the above classes are introduced in the translations. Dropped Content: A significant clause in the translation is removed. Please identify that error.  Source: Auf dieser Seite sind die Baudenkmäler der oberbayerischen Großen Kreisstadt Landsberg am Lech zusammengestellt.} \\
\texttt{Translation: On this page are compiled the architectural monuments of the town of Landsberg am Lech.} \\
\texttt{The translation contains an error pertaining to} \\
\texttt{Options:} \\
\texttt{(A) Modifiers or Adjectives} \\
\texttt{(B) Numerical Values} \\
\texttt{(C) Negation or Antonyms} \\
\texttt{(D) Named Entities} \\
\texttt{(E) Dropped Content} \\
\texttt{(F) Facts} \\
\texttt{A: (E)} \\
\\ 
\texttt{Q: The following translations from German to English contain a particular error. That error will be one of the following types: Named Entities: An entity (names, places, locations, etc.) is changed to a different entity. Numerical Values: Numerical values (ordinals or cardinals), dates, and/or units are changed. Modifiers or Adjectives: The modifiers and adjectives pertaining to a noun are changed. Negation or Antonyms: Introduce or remove a negation or change comparatives to their antonyms. Facts: Trivial factual errors not pertaining to the above classes are introduced in the translations. Dropped Content: A significant clause in the translation is removed. Please identify that error.  Source: Łeba ist eine Kleinstadt und ein Badeort im Powiat Lęborski der polnischen Woiwodschaft Pommern.} \\
\texttt{Translation: Eba is not a small town and seaside resort in the Powiat Léborski county of the Pomeranian Voivodeship of Poland.} \\
\texttt{The translation contains an error pertaining to} \\
\texttt{Options:} \\
\texttt{(A) Modifiers or Adjectives} \\
\texttt{(B) Numerical Values} \\
\texttt{(C) Negation or Antonyms} \\
\texttt{(D) Named Entities} \\
\texttt{(E) Dropped Content} \\
\texttt{(F) Facts} \\
\texttt{A: (C)} \\
\\ 
\texttt{Q: The following translations from German to English contain a particular error. That error will be one of the following types: Named Entities: An entity (names, places, locations, etc.) is changed to a different entity. Numerical Values: Numerical values (ordinals or cardinals), dates, and/or units are changed. Modifiers or Adjectives: The modifiers and adjectives pertaining to a noun are changed. Negation or Antonyms: Introduce or remove a negation or change comparatives to their antonyms. Facts: Trivial factual errors not pertaining to the above classes are introduced in the translations. Dropped Content: A significant clause in the translation is removed. Please identify that error.  Source: Georg Otto Bussenius war ein deutscher Lehrer.} \\
\texttt{Translation: Otto was a German teacher.} \\
\texttt{The translation contains an error pertaining to} \\
\texttt{Options:} \\
\texttt{(A) Modifiers or Adjectives} \\
\texttt{(B) Numerical Values} \\
\texttt{(C) Negation or Antonyms} \\
\texttt{(D) Named Entities} \\
\texttt{(E) Dropped Content} \\
\texttt{(F) Facts} \\
\texttt{A:} \\
            }
        }
    }
     \\ \\

    \theutterance \stepcounter{utterance}  
    & \multicolumn{4}{p{0.6\linewidth}}{
        \cellcolor[rgb]{0.561,0.737,0.561}{
            \makecell[{{p{\linewidth}}}]{
                \texttt{\tiny{[P1$\rangle$GM]}}
                \texttt{Based on the provided plans and previous attempts, I will carefully read the options and consider each type of error before making a decision. } \\
\\ 
\texttt{The source sentence is: Georg Otto Bussenius war ein deutscher Lehrer.} \\
\texttt{The translation is: Otto was a German teacher.} \\
\\ 
\texttt{To identify the type of error, I will examine the changes in the translation:} \\
\\ 
\texttt{1. Named Entities: The source sentence mentions "Georg Otto Bussenius," which is a specific person's name. The translation only mentions "Otto," which is a shortened version of the name. This could be considered a change in a named entity.} \\
\texttt{2. Numerical Values: There are no numerical values in the source sentence or the translation.} \\
\texttt{3. Modifiers or Adjectives: The source sentence uses the adjective "deutscher" to describe the teacher. The translation uses the adjective "German" to describe the teacher, which is a correct translation of the original adjective.} \\
\texttt{4. Negation or Antonyms: There is no negation or antonym in the source sentence or the translation.} \\
\texttt{5. Facts: There is no trivial factual error in the translation.} \\
\texttt{6. Dropped Content: There is no significant clause removed in the translation.} \\
\\ 
\texttt{Based on the analysis, the error in the translation is related to named entities, as the full name "Georg Otto Bussenius" is not translated correctly.} \\
\\ 
\texttt{The correct answer is: (D) Named Entities.} \\
            }
        }
    }
    & & \\ \\

    \theutterance \stepcounter{utterance}  
    & & & \multicolumn{2}{p{0.3\linewidth}}{
        \cellcolor[rgb]{0.9,0.9,0.9}{
            \makecell[{{p{\linewidth}}}]{
                \texttt{\tiny{[GM$|$GM]}}
                \texttt{Target: (D)} \\
            }
        }
    }
    & & \\ \\

    \theutterance \stepcounter{utterance}  
    & & & \multicolumn{2}{p{0.3\linewidth}}{
        \cellcolor[rgb]{0.9,0.9,0.9}{
            \makecell[{{p{\linewidth}}}]{
                \texttt{\tiny{[GM$|$GM]}}
                \texttt{game\_result = LOSE} \\
            }
        }
    }
    & & \\ \\

    \theutterance \stepcounter{utterance}  
    & & \multicolumn{4}{p{0.6\linewidth}}{
        \cellcolor[rgb]{0.95,0.95,0.95}{
            \makecell[{{p{\linewidth}}}]{
                \texttt{\tiny{[P1$\leftrightarrow$P]}}
                \texttt{Post{-}episode Reflexion (not sent to GM):} \\
\texttt{New plan: In my previous attempts, I failed to identify the correct type of error in the translations. Upon reflection, I realize that I was too focused on the changes in the translations and didn't carefully read the options. In the future, I will carefully read the options and consider each type of error before making a decision. Specifically, I will make sure to pay attention to the changes in named entities, numerical values, modifiers or adjectives, negation or antonyms, facts, and dropped content. If I'm unsure, I will take a closer look at the source sentence and the translation to determine the correct type of error.} \\
\\ 
\texttt{However, in this specific case, I failed to identify the correct type of error again. Upon re{-}examining the source sentence and the translation, I realize that I should have noticed that the translation "Otto was a German teacher" is missing the first name "Georg" which is a significant part of the original name "Georg Otto Bussenius". This is a clear example of "Dropped Content" error.} \\
\\ 
\texttt{In the future, I will make sure to pay closer attention to the content of the translation and compare it to the source sentence to identify any missing or dropped information.} \\
            }
        }
    }
    & \\ \\

\end{supertabular}
}

\end{document}
